\newpage
\section{Introdução}
\label{sec:revisao-introducao}

A dança de salão, e em particular o forró universitário, é uma prática de
ampla penetração cultural no Brasil, com benefícios documentados à saúde
física, à percepção corporal e à integração social de seus praticantes
\cite{havt2026dancar, fonseca2012, paiva2019danca}. Apesar de sua
popularidade, o aprendizado técnico da modalidade permanece dependente da
presença de um instrutor qualificado, capaz de observar e corrigir aspectos
como transferência de peso, alinhamento postural, conformidade rítmica e
execução de passos. A ausência de ferramentas computacionais acessíveis que
viabilizem o treinamento individual com feedback objetivo representa uma
lacuna relevante, especialmente em contextos onde o acesso a professores é
limitado por custo, disponibilidade ou localização geográfica.

O projeto Dança~AI propõe o desenvolvimento de um aplicativo Android em Kotlin
que utiliza visão computacional para preencher essa lacuna. A partir da câmera
do dispositivo e do framework MediaPipe \cite{lugaresi2019mediapipe,
bazarevsky2020blazeposeondevicerealtimebody}, o sistema analisa em tempo real
quatro dimensões do movimento do praticante: (i)~transferência de peso,
classificando o apoio como esquerdo, direito ou neutro a partir dos landmarks
dos membros inferiores; (ii)~postura corporal, avaliando o alinhamento do
tronco e dos ombros por regras geométricas e modelo de classificação embarcado;
(iii)~ritmo, por meio de metrônomo configurável com feedback auditivo e
validação pela comparação entre o BPM definido e a cadência detectada na
transferência de peso; e (iv)~identificação de movimentos, com classificador
embarcado para seis passos do forró (Balanço, Base~1, Base~2, Abertura, Giro
Invertido e Giro em Linha) nos modos Livre e Desafio. Toda a inferência é
executada localmente no dispositivo, sem dependência de conectividade
\cite{lee2019ondevice, li2020edgeai}.

Este documento está organizado da seguinte forma: a Seção~2 apresenta a
motivação do projeto; a Seção~3 sintetiza a revisão bibliográfica; a Seção~4
enumera os objetivos específicos; as Seções~5 e~6 descrevem os materiais,
métodos e resultados esperados; e a Seção~7 apresenta o cronograma de execução.
